\documentclass[aps,preprint]{revtex4}%
\usepackage{amsfonts}%
\usepackage{amsmath}%
\setcounter{MaxMatrixCols}{30}%
\usepackage{amssymb}%
\usepackage{graphicx}
%TCIDATA{OutputFilter=latex2.dll}
%TCIDATA{Version=5.50.0.2953}
%TCIDATA{CSTFile=revtex4.cst}
%TCIDATA{Created=Friday, May 30, 2008 08:28:10}
%TCIDATA{LastRevised=Friday, May 30, 2008 10:27:56}
%TCIDATA{<META NAME="GraphicsSave" CONTENT="32">}
%TCIDATA{<META NAME="SaveForMode" CONTENT="1">}
%TCIDATA{BibliographyScheme=Manual}
%TCIDATA{<META NAME="DocumentShell" CONTENT="Articles\SW\REVTeX 4">}
%BeginMSIPreambleData
\providecommand{\U}[1]{\protect\rule{.1in}{.1in}}
%EndMSIPreambleData
\newtheorem{theorem}{Theorem}
\newtheorem{acknowledgement}[theorem]{Acknowledgement}
\newtheorem{algorithm}[theorem]{Algorithm}
\newtheorem{axiom}[theorem]{Axiom}
\newtheorem{claim}[theorem]{Claim}
\newtheorem{conclusion}[theorem]{Conclusion}
\newtheorem{condition}[theorem]{Condition}
\newtheorem{conjecture}[theorem]{Conjecture}
\newtheorem{corollary}[theorem]{Corollary}
\newtheorem{criterion}[theorem]{Criterion}
\newtheorem{definition}[theorem]{Definition}
\newtheorem{example}[theorem]{Example}
\newtheorem{exercise}[theorem]{Exercise}
\newtheorem{lemma}[theorem]{Lemma}
\newtheorem{notation}[theorem]{Notation}
\newtheorem{problem}[theorem]{Problem}
\newtheorem{proposition}[theorem]{Proposition}
\newtheorem{remark}[theorem]{Remark}
\newtheorem{solution}[theorem]{Solution}
\newtheorem{summary}[theorem]{Summary}
\newenvironment{proof}[1][Proof]{\noindent\textbf{#1.} }{\ \rule{0.5em}{0.5em}}
\begin{document}
\preprint{HEP/123-qed}
\title[Short title for running header]{Long title that appears on the title page}
\author{First Author}
\affiliation{My Institution}
\author{Second Author}
\affiliation{My Institution}
\author{Third Author}
\affiliation{Other Institution}
\keywords{one two three}
\pacs{PACS number}

\begin{abstract}
Shell document for REV\TeX{} 4.

\end{abstract}
\volumeyear{year}
\volumenumber{number}
\issuenumber{number}
\eid{identifier}
\date[Date text]{date}
\received[Received text]{date}

\revised[Revised text]{date}

\accepted[Accepted text]{date}

\published[Published text]{date}

\startpage{101}
\endpage{102}
\maketitle
\tableofcontents


A OET is \emph{distinct} from ones based on FORDO's, as different sets of
GSO's (with the same occupation numbers) can produce the \emph{same} FORDO
while forming different $N$-electron states, with different energies.  This
difference can be understood by considering the action of the subgroup
$U\left(  \boldsymbol{\theta}\right)  $ of the unitary group $U\left(
\mathcal{H}^{1}\right)  $ of unitary transformations of one electron space,
$\mathcal{H}^{1}$, to itself, defined by
\begin{align}
u\left(  \boldsymbol{\theta}\right)  \left(  \left\vert \psi_{1}\right\rangle
,\cdots,\left\vert \psi_{r}\right\rangle \right)   &  =\left(  \left\vert
\psi_{1}\right\rangle ,\cdots,\left\vert \psi_{r}\right\rangle \right)
\left(
\begin{array}
[c]{ccc}%
e^{i\theta_{1}} & \cdots & 0\\
\vdots & \ddots & \vdots\\
0 & \cdots & e^{i\theta_{r}}%
\end{array}
\right)  \nonumber\\
&  =\left(  e^{i\theta_{1}}\left\vert \psi_{1}\right\rangle ,\cdots
,e^{i\theta_{r}}\left\vert \psi_{r}\right\rangle \right)
\end{align}
first on the FORDO $D^{1}$
\begin{equation}
D^{1}=\sum_{1\leq j\leq r}N_{j}\left\vert \psi_{j}\right\rangle \left\langle
\psi_{j}\right\vert ,
\end{equation}
then on the $N$-electron state density operator $D^{N}$
\begin{equation}
D^{N}=\sum_{\substack{1\leq j_{1}<\,\cdots<j_{N}\leq r\\1\leq k_{1}%
<\,\cdots<k_{N}\leq r}}\bar{c}_{j_{1}\cdots j_{N}}c_{k_{1}\cdots k_{N}%
}\left\vert \psi_{j_{1}}\cdots\psi_{j_{N}}\right\rangle \left\langle
\psi_{k_{1}}\cdots\psi_{k_{N}}\right\vert .
\end{equation}
that contracts to $D^{1}$, noting that $D^{N}$ is not in general diagonal in
the basis $\left\{  \left.  \left\vert \psi_{j_{1}}\cdots\psi_{j_{N}%
}\right\rangle \right\vert 1\leq j_{1}<\,\cdots<j_{N}\leq r\right\}  $. The
subgroup $U\left(  \boldsymbol{\theta}\right)  $ leaves $D^{1}$ invariant as
\begin{equation}
u\left(  \boldsymbol{\theta}\right)  D^{1}u\left(  \boldsymbol{\theta}\right)
^{\dagger}=\sum_{1\leq j\leq r}N_{j}e^{i\theta_{j}}\left\vert \psi
_{j}\right\rangle \left\langle \psi_{j}\right\vert e^{-i\theta_{j}}%
=\sum_{1\leq j\leq r}N_{j}\left\vert \psi_{j}\right\rangle \left\langle
\psi_{j}\right\vert ,
\end{equation}
but it does not leave the state $D^{N}$ invariant,
\begin{align}
u\left(  \boldsymbol{\theta}\right)  D^{N}u\left(  \boldsymbol{\theta}\right)
^{\dagger}  & =\sum_{\substack{1\leq j_{1}<\,\cdots<j_{N}\leq r\\1\leq
k_{1}<\,\cdots<k_{N}\leq r}}\bar{c}_{j_{1}\cdots j_{N}}c_{k_{1}\cdots k_{N}%
}e^{i\left\{  \theta_{j_{1}}+\cdots+\theta_{j_{N}}\right\}  }e^{-i\left\{
\theta_{k_{1}}+\cdots+\theta_{k_{N}}\right\}  }\left\vert \psi_{j_{1}}%
\cdots\psi_{j_{N}}\right\rangle \left\langle \psi_{k_{1}}\cdots\psi_{k_{N}%
}\right\vert \nonumber\\
& \neq D^{N},
\end{align}
where $\left\{  \left.  c_{j_{1}\cdots j_{N}}\right\vert 1\leq j_{1}%
<\,\cdots<j_{N}\leq r\right\}  $ are the CI coeffcients of the state vector
forming $D^{N}.$ The non invariance of the $N$-electron state density operator
will be shown by simple examples in the proceeding, where we show explicitly
that Second Order Reduced Density Operator's (SORDO's) of AGP\ states (thus
any higher order density operators) are not in general invariant to such transformations.

The observation that many sets of GSO's map into the same FORDO while
producing different $N$-electron states allows one to construct concrete
examples of DMF's (which are functionals of matrix representations of FORDO's)
through a Levy-Lieb \cite{Levy1979, Lieb1983}\ type of construction. Simple
examples of such functionals will also be displayed in the proceeding.


\end{document}